% Related Work draft for the TopoSteal extension (topology- and
% prediction-aware proactive work stealing). Drop into the new paper;
% NOT part of the original toposteal.tex (different paper, different scope).
%
% IMPORTANT: bibliographic details below (author lists, page numbers) are
% from research-pass notes during this project, not re-verified against
% the primary sources at draft time. Confirm every \bibitem against the
% actual paper before submission - do not cite off this file blind.

\section{Related Work}
\label{sec:related-work-ext}

Section~\ref{sec:related-work} covered the topology-aware work-stealing
literature TopoSteal itself builds on. This extension's contribution -
combining physical-topology-aware victim selection with proactive,
prediction-driven stealing - sits in a gap between two more recent lines
of work, neither of which combines both signals.

\subsection{Proactive and Predictive Load Balancing}

Chung et al.~\cite{chung2023reactive} extend reactive load balancing with
an online, ML-regression-based proactive scheme: a dedicated monitoring
thread characterizes tasks and predicts per-process load one iteration
ahead, triggering offloading before imbalance manifests rather than after.
Evaluated on an adaptive-mesh-refinement PDE solver, it reports 1.5--3.4$\times$
speedups over reactive baselines. Critically, the authors identify
topology-awareness as unresolved future work, and describe (but do not
evaluate) a probability-based victim-selection scheme in an appendix -
the exact combination this extension implements and evaluates.

A2WS~\cite{a2ws2024} is architecturally closer: it extends a Chase-Lev
deque directly with \emph{preemptive} stealing - initiating a steal once a
process's ideal completion time estimate signals it will finish early,
rather than waiting for its deque to empty - and an analytical (non-ML)
victim-reweighting scheme, explicitly positioned against Chung et
al.'s approach as lower-overhead (no dedicated monitoring thread). A2WS
is explicitly topology-agnostic: its only locality notion is a logical
communication-radius parameter, not physical NUMA/cache distance, and it
assumes homogeneous tasks, naming per-task weighting as its own next step.

This extension's core claim is that combining what each of these papers
identifies as its own gap - physical topology-awareness (missing from
both) with proactive, prediction-driven stealing (present in both, but
never alongside real hardware distance) - measurably outperforms either
signal alone. That is the required three-way ablation in
Section~\ref{sec:eval-ablation}: topology-only reactive (TopoSteal
as published), prediction-only topology-agnostic (A2WS-shaped), and the
combined policy.

\subsection{Topology-Aware Work Stealing}

Beyond the prior work already cited in Section~\ref{sec:related-work}
(SLAW~\cite{guo2010}, Olivier et al.~\cite{olivier2012}, OpenStream~\cite{drebes2014}),
PufferFish~\cite{pufferfish2020} introduces Hierarchical Elastic Tasks that
grow or shrink across NUMA places depending on local imbalance, reporting
1.5--1.9$\times$ speedups over HClib and CilkPlus on memory-bound,
NUMA-sensitive benchmarks. Like the topology-aware lineage generally, it
is reactive: none of this line anticipates starvation ahead of time.

\subsection{Microsecond-Scale, Low-Tail-Latency Scheduling}

A separate, more heavily-resourced lineage - ZygOS~\cite{zygos2017},
Shenango~\cite{shenango2019}, Caladan~\cite{caladan2020}, and
Perséphone~\cite{persephone} - targets P99/P99.9 tail latency for
microsecond-scale RPC-style tasks via kernel-bypass runtimes with
dedicated core-allocation layers (an IOKernel or equivalent), reactive
per-core work stealing (ZygOS), fine-grained (5\,$\mu$s) core reallocation
(Shenango), and interference mitigation via multicast IPIs (Caladan).
This extension does not compete in that space - building a comparable
kernel-bypass OS layer is out of scope here - but borrows its evaluation
discipline: Section~\ref{sec:eval-ablation} reports P99/P99.9 completion
latency and jitter alongside throughput, not throughput alone, precisely
because this lineage established that average-case numbers hide the
behavior that matters for interactive and inference-serving workloads.

\subsection{Summary}

No prior work combines physical hardware-topology-aware victim selection
with proactive/anticipatory stealing in a single mechanism; two
independent, recent papers~\cite{chung2023reactive,a2ws2024} each name the
piece the other lacks as their own unaddressed future work. This
extension evaluates that combination directly, on an AI-inference-adjacent
workload (Section~\ref{sec:eval-workload}) chosen for its irregular,
memory-bandwidth-bound task shape rather than as a claimed AI contribution
in itself.

% --- New \bibitem entries for this section (append to the existing
%     bibliography; do not renumber/duplicate keys already in
%     toposteal.tex's \bibitem{blumofe1999} etc.) ---

\begin{thebibliography}{99}

\bibitem{chung2023reactive}
M.~T. Chung, J.~Weidendorfer, K.~F\"{u}rlinger, and D.~Kranzlm\"{u}ller,
``From reactive to proactive load balancing for task-based parallel
applications in distributed memory machines,''
\emph{Concurrency and Computation: Practice and Experience}, vol.~35,
no.~24, 2023.

\bibitem{a2ws2024}
``Adaptive Asynchronous Work-Stealing for distributed load-balancing in
heterogeneous systems,'' arXiv:2401.04494, 2024.
% TODO: confirm author list from the source before citing in the
% submitted paper - not re-verified here.

\bibitem{pufferfish2020}
V.~Kumar et al., ``PufferFish: NUMA-Aware Work-stealing Library using
Elastic Tasks,'' in \emph{Proc. IEEE International Conference on High
Performance Computing, Data, and Analytics (HiPC)}, 2020.

\bibitem{zygos2017}
G.~Prekas, M.~Kogias, and E.~Bugnion, ``ZygOS: Achieving Low Tail Latency
for Microsecond-scale Networked Tasks,'' in \emph{Proc. ACM Symposium on
Operating Systems Principles (SOSP)}, 2017.
% TODO: confirm author order before citing in the submitted paper.

\bibitem{shenango2019}
A.~Ousterhout, J.~Fried, J.~Behrens, A.~Belay, and H.~Balakrishnan,
``Shenango: Achieving High CPU Efficiency for Latency-sensitive Datacenter
Workloads,'' in \emph{Proc. USENIX Symposium on Networked Systems Design
and Implementation (NSDI)}, 2019.

\bibitem{caladan2020}
J.~Fried, Z.~Ruan, A.~Ousterhout, and A.~Belay, ``Caladan: Mitigating
Interference at Microsecond Timescales,'' in \emph{Proc. USENIX Symposium
on Operating Systems Design and Implementation (OSDI)}, 2020.

\bibitem{persephone}
``Pers\'{e}phone: A kernel-bypass OS scheduler for dynamic, unpredictable
service times at microsecond timescales.''
% TODO: confirm exact title, authors, and venue before citing - noted
% from a research pass, not re-verified against the primary source.

\end{thebibliography}
